%=============================================================================
%  chapter2.tex -- The Atoms
%
%  Uses macros from thesis_main.tex Part 2 only. No notation is defined here.
%
%  BEAT: we build the instrument that will test the obvious suspect. Establish
%  that the inputs are trustworthy, and be honest about where they are not.
%
%  Every \section carries a % QUESTION: comment (thesis_main.tex:19).
%=============================================================================

\chapter{The Atoms}
\label{ch:atoms}

Chapter~\ref{ch:intro} left us with a question and no instrument to answer it. A
spectroscopist measures the ratio of two hydrogen lines in a divertor, looks the
answer up in a table, and reads off a temperature and a density. The table was
computed on the assumption that the plasma has stopped changing; a divertor
never stops changing. Everyone asked whether the table still holds worries about
the same thing first: the excited states of the atom are the fast, delicate part
of the system, and surely they cannot keep up with a plasma torn about on
microsecond timescales. Testing that suspicion needs a model that follows every
level of the atom in time, assuming in advance that none of them has settled.

This chapter builds the raw material for that model. There is no time derivative
anywhere in it. It is about the \emph{atoms}: what can happen to a hydrogen atom
in a plasma, how fast, where the numbers come from, and how far they can be
trusted. Chapter~\ref{ch:theory} puts them in motion.

One sentence belongs before any physics, because it is the reader's fastest
route into everything that follows. \textbf{The approximation this thesis tests
is the Bodenstein steady-state approximation, applied to atomic energy levels
instead of radical concentrations.} Anyone who has written
$\mathrm{d}[\mathrm{R}^{\bullet}]/\mathrm{d}t \approx 0$ for a reactive
intermediate has used it. The 43 energy levels below are the chemical species,
the atomic processes are the elementary reactions, and this chapter measures the
rate constants before that approximation is made.

%-----------------------------------------------------------------------------
\section{What happens to a hydrogen atom in a plasma}
\label{sec:atom_processes}
% QUESTION: what can happen to a hydrogen atom, and on what timescale does each
% thing happen?

\subsection{The atom, from scratch}

A hydrogen atom is one electron bound to one proton. Quantum mechanics allows
the electron only certain energies, labelled by the \emph{principal quantum
number} $n = 1,2,3,\dots$, with binding energy
\begin{equation}
  I_n = \frac{I_{\mathrm H}}{n^{2}},
  \qquad I_{\mathrm H} = \SI{13.6058}{\electronvolt}.
  \label{eq:binding}
\end{equation}

\noindent $I_n$ (\si{\electronvolt}) is the energy needed to pull the electron
free and leave a bare proton. Both sides are an energy, so the dimensional check
is trivial; the \emph{shape} of Eq.~\eqref{eq:binding} is what matters. The
levels crowd together as $n$ grows: $n=1$ lies \SI{10.20}{\electronvolt} below
$n=2$, but $n=14$ and $n=15$ are separated by \SI{0.0089}{\electronvolt}. High
levels are easy to move between and easy to strip off; low levels are not. That
asymmetry generates the two timescales of Chapter~\ref{ch:theory}.

Within a shell the electron also carries \emph{orbital angular momentum},
quantised by $\ell = 0,1,\dots,n-1$. In hydrogen, and only in hydrogen, all
$\ell$ states of a given $n$ have the same energy: they are \emph{degenerate}.
Spectroscopy labels them $\ell = 0,1,2,3,\dots \equiv
\mathrm{S},\mathrm{P},\mathrm{D},\mathrm{F},\dots$, so ``$2p$'' means $n=2$,
$\ell=1$. The number of quantum states sharing a label is the
\emph{statistical weight}
\begin{equation}
  g_{n\ell} = 2(2\ell+1),
  \qquad g_n = \sum_{\ell=0}^{n-1} g_{n\ell} = 2n^{2},
  \label{eq:weights}
\end{equation}

\noindent the factor 2 counting electron spin. $g$ is a dimensionless counting
number and appears only where a thermodynamic argument needs to know how many
ways a state can be occupied: $g = 6$ for $2p$, $g = 2$ for $2s$.

An atom dropping to a lower level can emit the energy difference as one photon.
The levels are discrete, so the photon energy is too, and the light appears at
one sharp wavelength: a \emph{spectral line}. Lines ending on $n=2$ form the
Balmer series and lie in the visible, which is why a fibre-coupled spectrometer
sees them; this thesis uses $H_\alpha$ ($n=3\to2$, \SI{656}{\nano\metre}) and
$H_\beta$ ($n=4\to2$, \SI{486}{\nano\metre}). A line's brightness is
proportional to the population of the level it comes from, so measuring light
measures populations, which is why the populations are worth computing.

\subsection{Six processes, and how fast they run}

Six things can happen to an atom in a plasma of electron density $\nel$
(\si{\per\cubic\centi\metre}) and temperature $\Te$ (\si{\electronvolt}), and
the model is built from these and nothing else. \emph{Electron-impact
excitation} pushes the bound electron up, and needs a free electron carrying at
least the energy gap. \emph{Electron-impact de-excitation} is the reverse and
has no threshold, so it is generally the faster of the pair. \emph{Ionisation}
removes the electron altogether, and is the only process that takes an atom out
of the set of levels we track. \emph{Recombination} does the reverse, by two
distinct mechanisms treated in Section~\ref{sec:ion_recomb}.
\emph{Spontaneous emission} drops the atom to a lower level with no collision
partner at all, the only process whose rate is independent of density.
\emph{Proton-impact $\ell$-mixing} shuffles population between the $\ell$
sublevels of one shell; it costs no energy, because those sublevels are
degenerate, and Section~\ref{sec:lmixing} shows it is the fastest process in the
model.

The chemical-engineering reading is exact. Each level is a species and each
process an elementary reaction. Excitation, de-excitation, ionisation and
$\ell$-mixing need a collision partner, so their rates carry a density factor
exactly as a second-order rate law does. Spontaneous emission is unimolecular,
like a first-order decomposition. Recombination is the odd one out, because it
creates atoms out of a species we have chosen not to track, and
Chapter~\ref{ch:theory} shows it therefore behaves like a constant feed stream
into a stirred tank rather than a reaction inside it.

Before quoting numbers it is worth saying what we expect, so the table can
disagree. An allowed optical transition should run at
$\sim10^{8}\,\si{\per\second}$; collisions at
$\nel \sim \SI{e14}{\per\cubic\centi\metre}$ should be slower, since a
coefficient of order $10^{-9}\,\si{\cubic\centi\metre\per\second}$ times that
density is only $10^{5}\,\si{\per\second}$; and ground-state ionisation at a few
electronvolts should be slower still, needing an electron above
\SI{13.6}{\electronvolt} from a distribution whose mean energy is far below it.
What $\ell$-mixing will do is not obvious in advance.

\begin{table}[tbp]
  \centering
  \caption[Process rates at the benchmark point]{Representative rates for the
  six processes at the benchmark point $\Te = \benchTe$, $\nel = \benchne$
  (grid indices $[23,5]$). Rate coefficients are read from the stored arrays of
  Table~\ref{tab:data_sources} at temperature index 23; a rate in
  \si{\per\second} is the coefficient times $\nel$, except for spontaneous
  emission, which needs no partner. The recombination entry is a feed per unit
  proton density, not a rate an atom experiences. The spread is ten orders
  of magnitude, which is why the problem is numerically stiff.}
  \label{tab:process_rates}
  \begin{tabular}{llrr}
    \toprule
    Process & Example & Rate (\si{\per\second}) & $1/\text{rate}$ \\
    \midrule
    Proton $\ell$-mixing      & $2s \to 2p$   & $1.475\times10^{11}$ & \SI{6.78}{\pico\second} \\
    Spontaneous emission      & $2p \to 1s$   & $6.268\times10^{8}$  & \SI{1.60}{\nano\second} \\
    Ionisation                & $2p \to$ H$^+$& $8.639\times10^{6}$  & \SI{116}{\nano\second} \\
    Collisional de-excitation & $2p \to 1s$   & $8.173\times10^{5}$  & \SI{1.22}{\micro\second} \\
    Collisional excitation    & $1s \to 2p$   & $7.686\times10^{4}$  & \SI{13.0}{\micro\second} \\
    Ionisation                & $1s \to$ H$^+$& $1.409\times10^{4}$  & \SI{71.0}{\micro\second} \\
    Recombination feed        & H$^+ \to 1s$  & $1.095\times10^{1}$  & n/a \\
    \bottomrule
  \end{tabular}
\end{table}

Table~\ref{tab:process_rates} confirms the expectations and exceeds one of them
by three orders of magnitude. Radiative decay of $2p$ takes
\SI{1.60}{\nano\second} and ground-state ionisation \SI{71.0}{\micro\second},
both as anticipated, but $\ell$-mixing empties $2s$ in under
\SI{7}{\pico\second}, two hundred times faster than the fastest radiative decay
in the atom. That number is why Section~\ref{sec:lmixing} exists. It is also the
first sight of the difficulty: a chemical engineer recognises a rate spread of
ten decades as a stiff network, of exactly the kind that motivates the
steady-state approximation for intermediates in the first place.

%-----------------------------------------------------------------------------
\section{The state space}
\label{sec:state_space}
% QUESTION: which energy levels does the model track, and why is l resolved
% below n=9 but not above it?

Hydrogen has infinitely many bound levels, so a model must stop somewhere, and
must also decide how finely to describe each level it keeps. Too few levels and
the cascade feeding the Balmer lines is cut short; too much detail at high $n$
and the matrix grows without the extra states carrying independent information.

This model tracks \textbf{43 states}, drawn in Figure~\ref{fig:state_space}
together with the five processes that connect them.
Levels $n=1$--$8$ are resolved in $\ell$
(indices 0--35, ordered $1S$, $2S$, $2P$, $3S$, \dots, $8J$); levels $n=9$--$15$
are \emph{bundled}, one state per shell carrying the full shell weight
$g_n = 2n^{2}$ (indices 36--42). The ordering is fixed by
\texttt{data/processed/collisions/K\_exc\_full/state\_index.csv} and every array
in the pipeline is indexed in it. The ion H$^{+}$ is deliberately \emph{not} a
44th state: it is an external reservoir, a choice whose consequences are listed
in Section~\ref{sec:atoms_scope}.

\begin{figure}[tbp]
  \centering
  \includegraphics[width=\textwidth]{fig2_1_state_space}
  \caption{\CapFigStateSpace}
  \label{fig:state_space}
\end{figure}

\paragraph{Why $\ell$ must be resolved at low $n$.} The obvious guess is that
$\ell$ is a refinement, nice to have, not load-bearing. At $n=2$ it is wrong,
and a selection rule is the reason. Electric-dipole (E1) radiation, much the
strongest emission mechanism, requires $\Delta\ell = \pm1$. The transition
$2p \to 1s$ satisfies this and runs at
$A_{2p\to1s} = \SI{6.2684e8}{\per\second}$
(\texttt{Radiative/A\_resolved.npy}, element $[0,2]$). The transition
$2s \to 1s$ does not: both states have $\ell = 0$, and the E1 rate is exactly
zero. In this dataset the total radiative loss from $2s$ is $\gamma_{2s} = 0$
identically; the two-photon decay hydrogen really possesses, of order
\SI{10}{\per\second}, is not included, and would be negligible beside the
collisional rates of Table~\ref{tab:process_rates} in any case. So the two
sublevels of $n=2$ are not refinements of each other: one empties in
\SI{1.6}{\nano\second} and the other cannot empty radiatively at all. Bundling
them and weighting the decay statistically would hand the metastable $2s$ level
a decay rate it does not have.

\paragraph{Why bundling is safe at high $n$.} Above $n=8$ the argument reverses.
Proton $\ell$-mixing scales roughly as $n^{4}$ (Section~\ref{sec:lmixing}), so
by $n \approx 9$ it redistributes population within a shell far faster than
anything else happens to that shell. The sublevels then sit permanently in the
ratio of their statistical weights and one number per shell holds all the
information.

That is an argument, not a measurement, and it can be tested: do the $\ell$
sublevels of a shell the model \emph{does} resolve come out statistical anyway?
For $n=4$ the statistical share of the $4f$ sublevel is $g_{4f}/g_4 = 14/32 =
0.4375$. Solving the model across the grid gives $0.4361$--$0.4375$ (register
entry A6), statistical to four figures. The largest departure,
\SI{0.3}{\percent}, falls at $\nel = \SI{e12}{\per\cubic\centi\metre}$, the
lowest density on the grid and exactly where $\ell$-mixing is weakest. Two caveats
belong in the same breath. The bundled block carries no $\ell$-mixing of its own:
it is \emph{assumed} statistical rather than driven to it, and that assumption
has not been tested directly. And the truncation at $n_{\max} = 15$ is a real
limitation whose size Chapter~\ref{ch:limits} reports as measured rather than
assumed.

\paragraph{Where the ordering and the grids come from.} A 43-state model has one
correct row ordering, and every array must agree on it; getting it wrong
produces no error, only plausible numbers about the wrong states. The pipeline
therefore never lets a script declare the ordering for itself.
\texttt{src/validation/cr\_context.py} loads the grids, the matrix and the state
ordering from the files the pipeline wrote, cross-checks their shapes, and
raises if any is missing or inconsistent; its docstring names the failure mode
it prevents: ``If a verification script hardcodes them and the pipeline later
changes, the script keeps running and silently mislabels every result'', which
is ``exactly the failure mode verification is supposed to catch''. Every
verification result in Chapters~\ref{ch:validation} and~\ref{ch:results} comes
from a script that imports it. The grid itself is declared in
\texttt{src/rates/assemble\_cr\_matrix.py} and written to disk so downstream code
reads rather than redeclares it: 50 temperatures logarithmically spaced over
$\Te = 1$--\SI{10}{\electronvolt} and 8 densities over
$\nel = 10^{12}$--\SI{e15}{\per\cubic\centi\metre}, giving 400 points. The
benchmark point carried through this thesis is the one nearest
$\Te = \SI{3}{\electronvolt}$, $\nel = \SI{e14}{\per\cubic\centi\metre}$: index
$[23,5]$, $\Te = \benchTe$, $\nel = \benchne$. It is an illustrative anchor, so
that the reader accumulates familiarity with one place in parameter space. No
conclusion rests on it; every general claim is carried by the full grid.

%-----------------------------------------------------------------------------
\section{From a cross section to a rate coefficient}
\label{sec:maxwellian}
% QUESTION: how does a cross section sigma(E), which is what atomic physics
% computes, become a rate constant a kinetic model can use?

Atomic collision theory does not deliver rate constants; it delivers
\emph{cross sections}. The cross section $\sigma(E)$ (\si{\square\centi\metre})
is the effective target area an atom presents to an electron of kinetic energy
$E$: a beam of flux $\Gamma$ (\si{\per\square\centi\metre\per\second}) drives the
process $\Gamma\sigma$ times per second. For excitation $\sigma$ vanishes below
the threshold $\Delta E$, rises above it, and falls again at high energy as the
electron passes too quickly to interact.

A plasma is not a beam. What a kinetic model needs is one number per
temperature, the \emph{rate coefficient} $C(\Te) = \langle\sigma v\rangle$,
which for a Maxwell--Boltzmann electron distribution reduces to
\begin{equation}
  C(\Te) = \sqrt{\frac{8}{\pi m_e}}\;(k_B\Te)^{-3/2}
           \int_{\Delta E}^{\infty}
           \sigma(E)\,E\,\mathrm{e}^{-E/k_B\Te}\,\mathrm{d}E,
  \label{eq:maxwell_average}
\end{equation}

\noindent with $m_e$ the electron mass (\si{\kilogram}), $k_B\Te$ the
temperature as an energy (\si{\joule}), $E$ the incident energy (\si{\joule}),
$\Delta E$ the threshold (\si{\joule}) and $\sigma$ in \si{\square\metre}; the
result is converted to \si{\cubic\centi\metre\per\second}. Dimensionally
$\si{\square\centi\metre}\times\si{\centi\metre\per\second} =
\si{\cubic\centi\metre\per\second}$, and $C\nel n_p$ is then
\si{\per\cubic\centi\metre\per\second}, a number of events per unit volume per
unit time: this is mass-action kinetics with $C$ as a second-order rate
constant. What Eq.~\eqref{eq:maxwell_average} does physically is weight the
cross section by how many electrons arrive at each energy and how fast, with the
lower limit enforcing that electrons below threshold contribute nothing.

One reading of it is a common trap. The prefactor carries $\Te^{-3/2}$, but
$C(\Te)$ does not fall as $\Te^{-3/2}$: the integral grows as more of the
distribution clears the threshold, and for a threshold process that growth wins
easily. The ground-state excitation coefficient $C_{1s\to2p}$ is
\SI{5.5316e-10}{\cubic\centi\metre\per\second} at the benchmark temperature and
nearly three orders of magnitude smaller at \SI{1}{\electronvolt}.

The implementation is \texttt{compute\_K\_CCC.py}, function
\texttt{maxwell\_average}, and four choices in it could have gone wrong
silently. The integral is trapezoidal on a 5000-point uniform grid; a 500-point
grid was tested and gives $\sim\SI{17}{\percent}$ error at
\SI{1}{\electronvolt}, so the grid density is measured, not defaulted. The
threshold is taken from Eq.~\eqref{eq:binding} rather than from the data file,
avoiding the noisy near-threshold region. Interpolation returns zero outside the
tabulated range. And the integral stops at the top of the data
(\SIrange{100}{968}{\electronvolt}) rather than at infinity; since
$\mathrm{e}^{-E/k_B\Te} = 4.5\times10^{-5}$ at \SI{100}{\electronvolt} and
$\Te = \SI{10}{\electronvolt}$, the missing tail is below \SI{0.01}{\percent}
over the range used here.

\subsection{Detailed balance, and why de-excitation is never computed directly}

Excitation and de-excitation are not independent. This is the statement a
chemical engineer makes when insisting that forward and reverse rate constants
of an elementary step be consistent with the equilibrium constant, and the
argument has the same shape. Put the plasma in complete thermodynamic
equilibrium, where every process is balanced by its reverse. For levels 1
(lower) and 2 (upper),
\begin{equation}
  n_1\,\nel\,C_{12} = n_2\,\nel\,C_{21}
  \quad\Longrightarrow\quad
  \frac{C_{12}}{C_{21}} = \left.\frac{n_2}{n_1}\right|_{\mathrm{eq}}
  = \frac{g_2}{g_1}\,\exp\!\left(-\frac{\Delta E}{k_B\Te}\right),
  \label{eq:detailed_balance}
\end{equation}

\noindent with $C_{12},C_{21}$ in \si{\cubic\centi\metre\per\second}, $g_1,g_2$
from Eq.~\eqref{eq:weights}, and $\Delta E = I_1 - I_2 > 0$. Both sides are
dimensionless. Physically: going up needs a rare fast electron and coming down
does not, so de-excitation wins by the Boltzmann factor, partly offset by there
being more upper states to fall out of.

The electron density cancels, because the same electron is the partner in both
directions. So do the populations: $C_{12}$ and $C_{21}$ are Maxwellian averages
of cross sections, which depend on $\Te$ and atomic structure and nothing else.
Equilibrium was only a device to fix the ratio; the ratio holds in any
Maxwellian plasma, however far the populations are from equilibrium. That is the
entire justification for using Eq.~\eqref{eq:detailed_balance} in a model whose
subject is departure from equilibrium. It fails only if the electrons themselves
are non-Maxwellian.

\paragraph{Worked number.} For $1s \to 2p$ at the benchmark temperature,
$g_1 = 2$, $g_2 = 6$, $\Delta E = \SI{10.20435}{\electronvolt}$, so
Eq.~\eqref{eq:detailed_balance} predicts $3\exp(-10.20435/2.947052) =
9.404800\times10^{-2}$. The stored arrays give $C_{1s\to2p}/C_{2p\to1s} =
\num{5.53155e-10}/\num{5.88163e-09} = 9.404800\times10^{-2}$, identical to seven
figures. Over all 819 transition pairs at that temperature, using $g$ and $I_n$
from \texttt{state\_index.csv}, the largest relative deviation is
$2.48\times10^{-9}$.

That agreement checks the implementation and nothing more, and must not be
reported as validating the atomic data. It is machine precision because it is
true by construction: de-excitation coefficients are never Maxwell-averaged from
de-excitation cross sections but derived from the excitation coefficients
through Eq.~\eqref{eq:detailed_balance}
(\texttt{compute\_K\_CCC.py}, function \texttt{detailed\_balance}). The reason
for building it that way is that it guarantees Saha--Boltzmann populations in
the high-density limit, which Chapter~\ref{ch:validation} uses as a gate, and
which is therefore a weaker gate than one the construction did not guarantee.
The independent statement, that the underlying cross sections themselves respect
microscopic reversibility, is a separate check reported in
Section~\ref{sec:atomic_uncertainty}.

\subsection{Two hydrogen energies, and why both are right}
\label{sec:two_energies}

A reader auditing the source will find two different numbers near
\SI{13.6}{\electronvolt} and should not treat the difference as sloppiness.
They are two distinct physical quantities.

The first is the \emph{Rydberg energy} $R_\infty hc =
\SI{13.605693122994}{\electronvolt}$, the ionisation energy of an idealised
hydrogen atom with an infinitely heavy, point-like, non-relativistic nucleus. It
is a defined constant of atomic theory, and it is the unit in which collision
calculations tabulate their energies. The second is the \emph{ionisation energy
of real hydrogen}, \SI{13.598434599702}{\electronvolt}, the measured energy
needed to remove the electron from a real H atom in its ground state. The second
is smaller by \SI{0.0534}{\percent}.

The pipeline uses each where it belongs, and the rule is simple. Cross sections
and thresholds supplied in Rydbergs must be converted back with the constant
they were tabulated in, so $R_\infty hc$ is carried by every collisional
quantity: the excitation thresholds of Section~\ref{sec:excitation}, the
ionisation thresholds of Section~\ref{sec:ion_recomb}, and $I_{\mathrm H}$
inside the $\ell$-mixing formula Eq.~\eqref{eq:psm20}. Photon energies and
spectroscopic level energies describe real atoms and carry the real ionisation
energy. Both constants were checked against their published values and both are
correct for what they are.

Two warnings go with that. First, the difference is \emph{not} simply the
reduced-mass correction: scaling $R_\infty hc$ by $\mu/m_e$ reproduces the
ionisation energy for about five digits and then stops, and the residual
$\SI{1.47e-4}{\electronvolt}$ is the Dirac fine-structure and Lamb-shift
contribution. Writing one as the other times $m_e/m_{\mathrm H}$ is wrong at the
$10^{-5}$ level. Second, the Rydberg energy is stored in the source under a name,
\texttt{CHI\_H}, conventionally reserved for the ionisation potential.
Anything that reads that symbol expecting an ionisation potential, a Saha or
Boltzmann exponent in particular, is \SI{0.0534}{\percent} too large. Nothing in
the assembled matrix does so, but the name is a trap and is flagged here rather
than renamed.

%-----------------------------------------------------------------------------
\section{Electron-impact excitation}
\label{sec:excitation}
% QUESTION: where do the excitation cross sections come from, and what was
% deliberately left out of them?

Excitation lifts atoms out of the ground state and populates everything the
spectrometer sees, so its cross sections matter more than any other input. They
come from the \emph{convergent close-coupling} (CCC) method of Bray and
co-workers~\cite{BrayStelbovics1992,Bray2002}, in which the collision
problem is solved by
expanding the target in a large square-integrable basis and enlarging that basis
until the answer stops moving. The ``convergent'' is the claim that the limit
is reached rather than estimated. The dataset is the 3115 raw files supplied by
Prof.\ Igor Bray (personal communication, 22 February 2026) in
\texttt{data/raw/ccc/e-H\_XSEC\_LS/}, parsed by
\texttt{src/parsers/parse\_ccc.py}. The acquisition report for the same dataset
records 3117 files; 3115 are on disk, and the difference is noted here rather
than reconciled.

Three instructions came with the data and all three are load-bearing.

\paragraph{The filenames read backwards.} The convention is
\texttt{FINAL.INITIAL}, so \texttt{1S.2P} holds $2p \to 1s$
\emph{de}-excitation and the excitation $1s \to 2p$ lives in \texttt{2P.1S}.
Reading it the wrong way round would transpose the dataset with no obvious
symptom, so it is worth saying which direction the model uses: rate coefficients
are Maxwell-averaged from the \emph{excitation}-direction files, selected in
\texttt{compute\_K\_CCC.py} by the test $n_{\mathrm f} > n_{\mathrm i}$, and
de-excitation follows by Eq.~\eqref{eq:detailed_balance}. The
de-excitation-direction files are parsed and stored but never enter the rate
tables.

\paragraph{Transitions with $\Delta n = 0$ were excluded, on the data
provider's instruction.} The note reads: ``Transitions with
$n_{\mathrm{final}} = n_{\mathrm{initial}}$ (e.g.\ 2S$\leftrightarrow$2P,
3S$\leftrightarrow$3P) have unphysically large cross sections due to
non-relativistic degeneracy in the CCC calculation. All $\Delta n = 0$
transitions are excluded from the database.'' The reason is that in a strictly
non-relativistic treatment the $\ell$ sublevels of a shell are exactly
degenerate, the transition energy is zero, and the long-range dipole interaction
gives a divergent cross section. After the exclusion 2190 blocks remain, all
with $\Delta n \ne 0$: 1320 in the excitation direction and 870 in the
de-excitation direction. The Week-2 acquisition report for the same dataset
records 1740, which is 870 of each; the parse was subsequently extended to the
$n=9$ and $n=10$ final states noted below, and the difference is stated here
rather than reconciled. The stated accuracy for what remains is
$\sim$\SI{5}{\percent}, which the benchmark of
Section~\ref{sec:atomic_uncertainty} is able to test.

That exclusion has a consequence which is easy to miss and which
Section~\ref{sec:lmixing} turns into the central argument of this chapter. With
$\Delta n = 0$ electron impact removed there is no electron-driven route between
$2s$ and $2p$ at all. Together with the near-absence of a radiative exit from $2s$, it leaves
the $n=2$ manifold with almost nothing to do, unless another collision partner
is admitted.

\paragraph{Covering the shells CCC does not reach.} CCC data extend to $n=8$
resolved in $\ell$, with resolved data reaching $n=9$ and $n=10$. Transitions
involving higher bundled shells are filled from the semi-empirical formulae of
Vriens and Smeets~\cite{VriensSmeets1980}, coded in \texttt{compute\_K\_VS.py};
the implemented expressions are their Eqs.~(17)--(19) directly, not the
equivalent rearrangement of Hartgers \emph{et al.}~\cite{Hartgers2001} that the
file's header docstring describes. Their accuracy is $\sim$\SI{20}{\percent}
against $\sim$\SI{5}{\percent} for CCC, so the merge uses CCC wherever it
exists. Those formulae need one further input, the shell-to-shell absorption
oscillator strength, built from the same Einstein coefficients used in
Section~\ref{sec:radiative}, with a Kramers fallback where a value is missing
that is recorded per transition in the metadata.
Table~\ref{tab:exc_coverage} records which source supplies which block, because
``the excitation data are from CCC'' would misdescribe a table that is a quarter
Vriens--Smeets.

\begin{table}[tbp]
  \centering
  \caption[Provenance of the excitation rate coefficients]{Which source supplies
  which excitation transitions, from the coverage record in
  \texttt{src/rates/assemble\_K\_exc.py} and the per-transition attribution in
  \texttt{K\_exc\_meta.csv}. ``res'' is an $\ell$-resolved state ($n \le 8$),
  ``bund'' a bundled shell ($n = 9$--$15$). The collapse rows are $\ell$-resolved
  CCC data summed over final $\ell$ to give a shell total, which keeps CCC
  accuracy for the first two bundled shells instead of falling back to
  Vriens--Smeets. De-excitation coefficients are not listed because they are not
  independent data: each follows from its excitation partner by
  Eq.~\eqref{eq:detailed_balance}.}
  \label{tab:exc_coverage}
  \begin{tabular}{llrl}
    \toprule
    Block & Source & Pairs & Stated accuracy \\
    \midrule
    res $\to$ res, $n \le 8$      & CCC                        & 546 & $\sim$\SI{5}{\percent} \\
    res $\to$ $n=9$ bundled       & CCC, collapsed over $\ell$ &  36 & $\sim$\SI{5}{\percent} \\
    res $\to$ $n=10$ bundled      & CCC, collapsed over $\ell$ &  36 & $\sim$\SI{5}{\percent} \\
    res $\to$ $n=11$--$15$        & Vriens--Smeets             & 180 & $\sim$\SI{20}{\percent} \\
    bund $\leftrightarrow$ bund   & Vriens--Smeets             &  21 & $\sim$\SI{20}{\percent} \\
    \midrule
    total                         &                            & 819 & \\
    \bottomrule
  \end{tabular}
\end{table}

%-----------------------------------------------------------------------------
\section{Ionisation and recombination}
\label{sec:ion_recomb}
% QUESTION: what connects the 43 bound levels to the ion reservoir outside them,
% and why is recombination quadratic in density?

\subsection{Ionisation: the only true loss}

Ionisation is the only process that removes an atom from the bound manifold
altogether; everything else moves an atom from one tracked level to another. In
reactor terms it is the irreversible step consuming reactant to product, and the
rest is internal shuffling. The ionisation rate coefficient $Q_p$
(\si{\cubic\centi\metre\per\second}) is obtained by Maxwell-averaging the CCC
\emph{total ionisation cross section} (TICS) with
Eq.~\eqref{eq:maxwell_average}, the lower limit being the binding energy of
Eq.~\eqref{eq:binding} (\texttt{compute\_K\_TICS.py}). The threshold is
$\ell$-independent in hydrogen but the cross sections are not, so the
coefficients are $\ell$-resolved for $n \le 8$.

The stored coefficients are not uniform in origin and the metadata records this
state by state. The 36 resolved levels carry \texttt{CCC\_TICS}; the $n=9$
bundle carries \texttt{CCC\_TICS\_direct}, taken from the bundled
\texttt{TICS.9} file rather than assembled from sublevels; the six shells
$n=10$--$15$ carry \texttt{Lotz1968}, the semi-empirical formula of
Lotz~\cite{Lotz1968}, with a note recorded against every one of them:
``overestimates CCC by factor \textasciitilde4--8; negligible impact''
(\texttt{collisions/tics/K\_ion\_final\_meta.csv}). That is an honest flag, not
a resolved question, and Section~\ref{sec:atomic_uncertainty} returns to it,
because the overestimate does not stay confined to ionisation.

\paragraph{Worked number.} At the benchmark point
$Q_{1s} = \SI{1.0140e-10}{\cubic\centi\metre\per\second}$ and
$Q_{2p} = \SI{6.2171e-8}{\cubic\centi\metre\per\second}$
(\texttt{K\_ion\_final.npy}, temperature index 23): about 600 times easier to
ionise an atom already in $n=2$, which has to supply only
\SI{3.40}{\electronvolt} instead of \SI{13.61}{\electronvolt}. This step-ladder
character is why the excited levels matter to a quantity as apparently unrelated
as the total ionisation rate.

\subsection{Recombination: two mechanisms, two density dependences}

Recombination runs the other way, and telling its two mechanisms apart matters
because they scale differently with density. In \emph{radiative recombination}
the excess energy leaves as a photon: one electron and one proton are needed, so
events go as $\nel\nion$ and the coefficient $\alpha_{\mathrm{RR}}$ has units
\si{\cubic\centi\metre\per\second}. In \emph{three-body recombination} a second
free electron carries the energy away: three particles must meet, so events go
as $\nel^{2}\nion$ and $\alpha_{\mathrm{3BR}}$ has units
\si{\centi\metre\tothe{6}\per\second}. Nothing here is special to plasmas. It is
the difference between a bimolecular and a termolecular elementary step, with
the same consequence: the second mechanism takes over as density rises.

Radiative recombination is computed shell by shell from Johnson's three-term
Gaunt-factor expansion~\cite{Johnson1972}, itself a fit to the treatment of
Seaton~\cite{Seaton1959},
\begin{equation}
  \alpha_{\mathrm{RR}}(n,\Te)
  = D\,y^{3/2}\,\mathrm{e}^{y}
    \left[g_0(n)E_1(y) + g_1(n)E_2(y) + g_2(n)E_3(y)\right],
  \qquad y = \frac{I_n}{k_B\Te},
  \label{eq:alpha_RR}
\end{equation}

\noindent with $D = \SI{5.197e-14}{\cubic\centi\metre\per\second}$, $E_i$ the
exponential integrals and $g_0,g_1,g_2$ the Gaunt polynomials of Johnson's
Table~1 (\texttt{recombination\_rates.py}). $y$ and the bracket are
dimensionless and $D$ carries the units, so both sides are
\si{\cubic\centi\metre\per\second}. Physically, capture into a shell is easiest
so capture is favoured at low electron energy. Over the range used here
$\alpha_{\mathrm{RR}}(n,\Te)$ decreases monotonically with $\Te$ for every
level, with no interior maximum: Eq.~\eqref{eq:alpha_RR} gives
$\alpha \propto \Te^{-1/2}$ at high energy and drives $\alpha \to 0$ at low
energy, so the maximum lies below \SI{1}{\electronvolt} for the shells that
matter here. The shell total is split over
$\ell$ statistically, $\alpha_{\mathrm{RR}}(n,\ell) =
\alpha_{\mathrm{RR}}(n)(2\ell+1)/n^{2}$~\cite{MaoKaastra2016}, a step requiring
$\nel \gtrsim \SI{e10}{\per\cubic\centi\metre}$, two decades below our lowest
density.

One property of Eq.~\eqref{eq:alpha_RR} needed checking rather than assuming.
$g_1(n)$ is negative and grows with $n$, so the three-term sum can in principle
turn negative for high shells at low temperature; an earlier analysis in this
project raised that objection and recommended the single-term form instead. The
code keeps the three-term expansion and guards it with a clamp at zero, which
would silently replace a negative value rather than report it. Reading the
stored arrays directly, neither \texttt{alpha\_RR\_resolved.npy} nor
\texttt{alpha\_RR\_bundled.npy} contains an exact zero anywhere on the 50-point
temperature grid: the clamp never fires over the range this thesis uses, though
it would need revisiting for a colder grid or a longer ladder.

Three-body recombination is not computed independently at all. It comes from
ionisation by detailed balance~\cite{Griem1997}, exactly as de-excitation came
from excitation:
\begin{equation}
  \alpha_{\mathrm{3BR}}(n,\ell,\Te)
  = Q_{n\ell}(\Te)\,\frac{g_{n\ell}}{2}
    \left(\frac{h^{2}}{2\pi m_e k_B \Te}\right)^{3/2}
    \exp\!\left(\frac{I_n}{k_B\Te}\right),
  \label{eq:alpha_3BR}
\end{equation}

\noindent where $h$ is Planck's constant and the bracketed factor is the cube of
the electron thermal de Broglie wavelength, a volume, so the units read
$\si{\cubic\centi\metre\per\second}\times\si{\cubic\centi\metre} =
\si{\centi\metre\tothe{6}\per\second}$, as required.
Equation~\eqref{eq:alpha_3BR} imposes Saha equilibrium between a level and the
continuum: in complete equilibrium the ionisation flux $Q\,n_{n\ell}\nel$ and
the recombination flux $\alpha_{\mathrm{3BR}}\nel^{2}\nion$ must balance, and
rearranging Saha for $\nel\nion/n_{n\ell}$ gives it directly. Like detailed
balance for de-excitation, this guarantees the high-density limit rather than
testing it.

Assembled, the two mechanisms give the source used throughout the thesis,
$\Svec = \nel(\bm{\alpha}_{\mathrm{RR}} + \nel\bm{\alpha}_{\mathrm{3BR}})$,
whose units are $\si{\per\cubic\centi\metre} \times
\si{\cubic\centi\metre\per\second} = \si{\per\second}$, so that $\Svec\nion$ is
atoms created per unit volume per unit time. Chapter~\ref{ch:theory} explains
why this is a feed stream rather than an element of the rate matrix.

\paragraph{Worked numbers, and a surprise.} The obvious guess is that three-body
recombination is a high-density correction, negligible at
$\nel = \SI{e12}{\per\cubic\centi\metre}$. That is true only for the ground
state. Forming the three-body fraction
$\nel\alpha_{\mathrm{3BR}}/(\alpha_{\mathrm{RR}}+\nel\alpha_{\mathrm{3BR}})$
from the stored coefficients at temperature index 23 gives, for $1s$, below
$10^{-4}$ at $\nel = \SI{e12}{\per\cubic\centi\metre}$ and $0.0085$ at
\SI{e15}{\per\cubic\centi\metre}; but for $n=15$ it gives $0.9996$ at the
\emph{lower} density and above $0.9999$ at the higher. Capture into the ground
state is radiative essentially everywhere on the grid; capture into the top
shell is three-body essentially everywhere on the grid. Equation~\eqref{eq:alpha_3BR}
shows why: $\alpha_{\mathrm{3BR}}$ inherits an ionisation coefficient that is
enormous for a shell bound by \SI{0.060}{\electronvolt}, and the exponential
that suppresses it for the ground state is close to unity at the top of the
ladder. A loosely bound electron is easy to knock off and, by the same token,
easy to capture in a three-body encounter.

That has a consequence Section~\ref{sec:atomic_uncertainty} takes up: the
three-body coefficients for $n=10$--$15$ are built by
Eq.~\eqref{eq:alpha_3BR} from Lotz ionisation coefficients the metadata itself
calls too large by 4--8, and those are precisely the shells whose feed is
entirely three-body.

%-----------------------------------------------------------------------------
\section{Radiative decay}
\label{sec:radiative}
% QUESTION: how does an atom lose energy without a collision partner, and which
% transitions are forbidden to it?

An excited atom need not be disturbed to decay. The \emph{Einstein coefficient
for spontaneous emission} $A_{qp}$ (\si{\per\second}) is the probability per unit
time that an atom in upper level $q$ drops to lower level $p$, and the total
radiative loss from $q$ is $\gamma_q = \sum_{p<q}A_{qp}$, also in
\si{\per\second}, with $1/\gamma_q$ the radiative lifetime. This is a
first-order rate constant in the ordinary chemical sense, the one process with
no concentration dependence, so the only one whose relative importance rises as
the plasma thins.

Not every pair of levels is connected. E1 radiation requires $\Delta\ell = \pm1$;
transitions violating it are \emph{forbidden}, meaning slower by many orders of
magnitude rather than impossible. The dataset here contains E1 transitions only,
so a forbidden transition has $A = 0$ exactly. The related \emph{oscillator
strength}, a dimensionless measure of how strongly a transition couples to the
radiation field, feeds the Vriens--Smeets formulae of
Section~\ref{sec:excitation}.

Two consequences of the selection rule matter to this thesis. First,
$2s \to 1s$ is forbidden, so $\gamma_{2s} = 0$ against
$\gamma_{2p} = \SI{6.2684e8}{\per\second}$. The $n=2$ shell holds one level
that empties in \SI{1.6}{\nano\second} and one that cannot empty radiatively.
Second, $4f \to 2p$ is forbidden because $\ell$ would change by two, so the $4f$
sublevel holds \SI{43.75}{\percent} of the $n=4$ shell by weight and contributes
nothing to $H_\beta$, decaying only to $3d$. The $n=4$ \emph{shell} population
and the $H_\beta$ \emph{line} intensity are therefore not proportional, and
Chapter~\ref{ch:validation} measures how much that matters to the observable.
The three channels that do make $H_\alpha$ are $3s \to 2p$
($\SI{6.3172e6}{\per\second}$), $3p \to 2s$ ($\SI{2.2461e7}{\per\second}$) and
$3d \to 2p$ ($\SI{6.4688e7}{\per\second}$); the three that make $H_\beta$ are
$4s \to 2p$ ($\SI{2.5796e6}{\per\second}$), $4p \to 2s$
($\SI{9.6735e6}{\per\second}$) and $4d \to 2p$ ($\SI{2.0637e7}{\per\second}$).

The coefficients come from the exact hydrogenic code of
Hoang-Binh~\cite{HoangBinh2005}, distributed through the Computer Physics
Communications Program Library as ADUU~v1.0, which evaluates the radial
integrals analytically rather than fitting them~\cite{HoangBinh1990}. Its 1015
transitions for $n \le 15$ are stored in
\texttt{Radiative/H\_A\_E1\_LS\_n1\_15\_physical.csv}; every entry has
$|\Delta\ell| = 1$ and $n_{\text{upper}} > n_{\text{lower}}$, and
\texttt{radiative\_rates.py} sweeps every allowed combination on load and raises
if one is missing rather than assuming completeness. For the bundled shells
the $A$ values are averaged over $\ell$ with statistical weights, the same
assumption that justified bundling in Section~\ref{sec:state_space}, applied
consistently. The pattern to notice is that $\gamma$ falls steeply with $n$,
from \SI{6.27e8}{\per\second} for $2p$ to \SI{8.16e4}{\per\second} for the
$n=15$ bundle: radiative decay dominates at low $n$ and collisions at high $n$,
and the level where they cross is a measurable feature of the model that
Chapter~\ref{ch:validation} locates.

\paragraph{These are not the recommended values, and it matters to say so.}
An external check of the four Lyman-series coefficients finds them
\SIrange{0.056}{0.060}{\percent} \emph{above} the critically evaluated values
recommended by Wiese and Fuhr~\cite{WieseFuhr2009}: for Lyman-$\alpha$,
\SI{6.2684e8}{\per\second} stored against \SI{6.2649e8}{\per\second}
recommended. The offset is systematic, not scattered, and it has a cause. What
this code computes is the \emph{exact non-relativistic, infinite-nuclear-mass
hydrogenic} transition probability, reproduced to 1--4 parts in $10^{5}$; the
recommended values include the finite proton mass and the small relativistic and
quantum-electrodynamic corrections that the non-relativistic Schr\"{o}dinger
problem does not contain. Both sets are correct for what they are, and using the
exact hydrogenic set throughout is internally self-consistent, which is what a
rate matrix needs. \textbf{But this dataset must not be described as the
recommended or evaluated set, and no evaluated compilation should be cited as
its source.} It is the analytic hydrogenic solution, and that is what is written
here.

\paragraph{The largest silent error available in a rate matrix.} One number in
this section can be wrong by a third without anything complaining, and the trap
is worth spelling out because it is the exact mistake this chapter's
$\ell$-resolution is designed to avoid. The value
$A_{2p\to1s} = \SI{6.2684e8}{\per\second}$ is an \emph{$\ell$-resolved level
rate}: it is the decay probability of an atom that is in $2p$. The
corresponding \emph{$n$-resolved shell rate} for $n=2 \to n=1$ is smaller by
the fraction of the shell that can radiate at all,
\begin{equation}
  A_{n=2 \to n=1}
  = \frac{g_{2p}}{g_{2}}\,A_{2p\to1s}
  = \frac{6}{8}\times\SI{6.2684e8}{\per\second}
  = \SI{4.7013e8}{\per\second},
  \label{eq:shell_rate}
\end{equation}

\noindent the factor $3/4$ being the statistical weight of $2p$ within $n=2$,
and the missing quarter being the $2s$ sublevel that has no E1 route to the
ground state at all. Both sides are in \si{\per\second}. Placing the
$\ell$-resolved number on an $n$-resolved matrix element therefore overestimates
Lyman-$\alpha$ emission by \SI{33.3}{\percent}, and nothing in a conservation
check or a sign check would notice, because the error is in a magnitude and
magnitudes cancel in column sums (Section~\ref{sec:atoms_corrections}).

This model is $\ell$-resolved for $n \le 8$, so the $\ell$-resolved value is the
correct one to use here, but that was audited rather than assumed. Reading the
stored block directly: the element $A[1s,2s]$ is exactly zero, the element
$A[1s,2p]$ carries \SI{6.2684e8}{\per\second}, $\gamma_{2s} = 0$ and
$\gamma_{2p} = \SI{6.2684e8}{\per\second}$, and across all
$36\times36$ resolved elements there is not one non-zero entry with
$|\Delta\ell| \neq 1$. The shell-averaged value of Eq.~\eqref{eq:shell_rate}
appears nowhere in the resolved block. The bundled block $n = 9$--$15$ is
$n$-resolved by construction and carries statistically weighted shell rates
throughout, which is the consistent treatment for states that are always in
statistical $\ell$-equilibrium.

\subsection{The photons are assumed to leave}
\label{sec:radiative_transfer}

Writing $\gamma_q$ as a bare sum of Einstein coefficients contains an
assumption: that every emitted photon escapes and never returns. A plasma for
which this holds is \emph{optically thin}. If it is optically thick, a photon
emitted by one atom is absorbed by another, which returns to the upper level,
and the effective decay rate falls. The chemical-engineering analogue is a
reaction slowed because the product cannot leave and keeps re-adsorbing:
the kinetics are unchanged, the effective rate constant is not.

The size of the effect is set by the \emph{optical depth} $\tau_c$, the number of
absorption lengths a photon must cross to escape, and carried by the
\emph{population escape factor} $\Theta_P$, defined so that the effective
coefficient is $\Theta_P A$. For the resonance line the model is most exposed
to, Lyman-$\alpha$ ($2p \to 1s$), \texttt{src/analysis/escape\_factor.py}
evaluates the escape factor of Holstein~\cite{Holstein1947} in the formulation
of Behringer~\cite{Behringer1998} adopted as routine ADAS214 of the Atomic Data
and Analysis Structure (ADAS)~\cite{Summers2006}, by direct quadrature and with
no fitting formula:
\begin{equation}
  \Theta_P(\tau_c)
  = \frac{1}{\sqrt{\pi}}\int_{-\infty}^{\infty}
    \mathrm{e}^{-x^{2}}\,
    \exp\!\left(-\tau_c\,\mathrm{e}^{-x^{2}}\right)\mathrm{d}x,
  \qquad
  \tau_c = \ngs\,\sigma_0\,\frac{D}{2},
  \label{eq:escape}
\end{equation}

\noindent with $x$ the frequency offset in Doppler widths (dimensionless),
$\ngs$ the ground-state density (\si{\per\cubic\centi\metre}), $\sigma_0$ the
line-centre absorption cross section (\si{\square\centi\metre}) from the
Ladenburg relation with oscillator strength $f_{12} = 0.4162$ and
$\lambda_0 = \SI{121.567}{\nano\metre}$, and $D$ twice the escape path length
(\si{\centi\metre}). Both sides of the first expression are dimensionless.
It says: weight the escape probability by the fraction of photons emitted at
each frequency, and note that photons in the line wings get out when photons at
the line centre cannot.

\paragraph{The geometry is mislabelled, and the label matters.} The module
documents Eq.~\eqref{eq:escape} as applying to ``a homogeneous slab of
half-thickness $b$''. Read directly from the ADAS214 manual, equation 3.14.14 is
the \emph{isotropic} case (the escape probability at the centre of a sphere,
ADAS geometry g1) and not the slab case, g2. The distinction is not cosmetic
in the thick regime: a true slab requires the further angular average
$\Theta_{\rm slab}(\tau) = \int_0^1 \Theta_{\rm iso}(\tau/\mu)\,\mathrm{d}\mu$,
which is smaller by a factor approaching two, measured as $0.529$ at
$\tau_c = 1$ and $0.477$--$0.484$ at $\tau_c = 100$--$1000$. Everything else in
the module was checked against the primary source and is correct: the convention
$\tau_c = \kappa_0 b$ with $b$ the edge-to-centre distance is exactly ADAS's own,
the Ladenburg constant $\pi e^{2}/m_ec = \SI{2.65400885e-2}{\square\centi\metre\hertz}$
is right, the Doppler profile normalisation and thermal speed reproduce ADAS's
own tabulated widths, and the leading Holstein constant is 1 rather than the
1.6 sometimes quoted. \textbf{This is reported, not repaired.} Where the plasma
is thin, $\Theta_P \to 1$, the two geometries agree and nothing depends on the
choice, and Section~\ref{sec:atoms_scope} and Chapter~\ref{ch:theory} show that
this is the whole region carrying this thesis's quantitative results. Where the plasma
is thick, at the cold edge, the escape factor computed here is roughly twice too
large if a slab is the right picture, and Chapter~\ref{ch:limits} carries that
caveat rather than resolving it.

\paragraph{Two constants in this module are wrong, and both are immaterial.}
An external audit of the same file found the elementary charge entered as
$e = \SI{4.80326e-10}{esu}$ against the CODATA value
$\num{4.803204713e-10}$~\cite{CODATA2018}, high by \SI{0.00115}{\percent} and
entering $\sigma_0$ as $e^{2}$, so \SI{0.0023}{\percent}; and the Doppler width
computed with $m = \SI{1.67262e-24}{\gram}$, which is the \emph{proton} mass
where the physically correct object is the neutral hydrogen \emph{atom},
\SI{1.6735328407e-24}{\gram}. A neutral atom Doppler-broadens with its own mass,
so the line is \SI{0.027}{\percent} too wide. Neither shifts any conclusion in
this thesis by a detectable amount, and neither has been edited: the repository
is under verification rather than development, and a constant corrected quietly
is a constant nobody can find again. They are recorded here so that a reader who
recomputes $\sigma_0$ and gets a slightly different number knows why.

\textbf{Every rate matrix in this thesis is built with bare $A$ coefficients,
$\Theta_P \equiv 1$.} Equation~\eqref{eq:escape} is a diagnostic applied
afterwards, to find where that choice is defensible. Chapter~\ref{ch:theory}
carries out that evaluation across the grid and Chapter~\ref{ch:limits} states
the consequences where the plasma is not thin. The question can be deferred
because $\tau_c$ depends on $\ngs$, an output of the model rather than an input,
and on a path length $D$ the model does not contain at all.

%-----------------------------------------------------------------------------
\section{Proton-impact \texorpdfstring{$\ell$}{l}-mixing, and why it is not optional}
\label{sec:lmixing}
% QUESTION: why include a collision process that most collisional-radiative
% treatments of hydrogen leave out?

\subsection{The physical picture, before any formula}

Collisional--radiative treatments of hydrogen often omit proton collisions on
the grounds that protons are slow and heavy and therefore ineffective. For
transitions that change $n$ that is right. For transitions that only change
$\ell$ it is exactly backwards, and the chain is worth following link by link,
because it explains a rate that would otherwise be merely a large number in a
table.

\emph{Degeneracy makes the mixing free, and slowness helps.} Within a shell,
$ns$, $np$, $nd$ have identical energy, so moving population among them costs
nothing: there is no threshold to clear, unlike excitation $n \to n+1$. And a
proton, 1836 times heavier than an electron, moves 43 times more slowly at the
same temperature, so the encounter lasts far longer. For a process needing a
sustained push rather than a sharp kick, a long gentle encounter is worth more.

\emph{The mechanism is the linear Stark effect.} A passing proton is a moving
charge; while nearby, the atom sits in a nearly static electric field. Because
$s$, $p$ and $d$ are degenerate, hydrogen responds linearly, and the field's own
eigenstates are mixtures of $s$, $p$ and $d$. An atom in a pure $np$ state is
therefore no longer in an eigenstate once the field arrives; it evolves as a
superposition, and when the proton leaves the population has been redistributed
across $\ell$. No photon and no dipole selection rule is involved.

\emph{Distant collisions dominate, logarithmically, and Debye screening cuts
them off.} The field at impact parameter $b$ falls as $1/b^{2}$, so the
transition probability does too; but the number of protons passing near $b$
grows as $b\,\mathrm{d}b$. The product is $\mathrm{d}b/b$, so every decade in
$b$ contributes equally and the integral produces a logarithm, which is why
the process is strong, and why it needs a cut-off. That cut-off is physical: a
proton is dressed by a screening cloud, its potential falls as
$\mathrm{e}^{-R/\lambda_D}/R$, and beyond about one Debye length $\lambda_D$ the
atom feels nothing. So $b_{\max}\approx\lambda_D$ and the Coulomb logarithm is
$\ln(\lambda_D/b_{\min})$.

Assembling those links by an impact-parameter calculation gives the scaling
\begin{equation}
  \nu_{\ell\text{-mix}} \;\propto\; \nion\, n^{4}\, T_i^{-1/2}\, \ln\Lambda,
  \label{eq:lmix_scaling}
\end{equation}

\noindent with $\nu$ in \si{\per\second}, $\nion$ the proton density
(\si{\per\cubic\centi\metre}), $T_i$ the \emph{proton} temperature
(\si{\electronvolt}) and $\ln\Lambda$ the Coulomb logarithm. The $n^{4}$ comes
from a dipole matrix element scaling as $n^{2}$ and entering squared; the
$T_i^{-1/2}$ from slower protons interacting for longer.
Equation~\eqref{eq:lmix_scaling} tells us in advance that $\ell$-mixing will be
fast, will get faster with density, and will grow explosively up the ladder,
which is the quantitative content of the bundling argument in
Section~\ref{sec:state_space}.

\subsection{The implemented formula}

The rates actually used come from the Debye-cutoff form of the Pengelly--Seaton
treatment~\cite{PengellySeaton1964} given by Badnell \emph{et
al.}~\cite{Badnell2021}, called PSM20 in the code and in this thesis, coded
in \texttt{src/rates/compute\_lmix.py}. For the downward transition within shell
$n$, writing $\ell_> = \max(\ell,\ell')$,
\begin{equation}
  q_{ji} = \frac{a_0^{3}}{\tau_0}
           \left(\frac{\pi\mu I_{\mathrm H}}{k_B\Te}\right)^{1/2}
           \frac{D_{ji}}{\omega_\ell}\,F(U_m),
  \qquad
  D_{ji} = 6n^{2}\ell_>\left(n^{2}-\ell_>^{2}\right),
  \label{eq:psm20}
\end{equation}

\noindent with $a_0$ the Bohr radius, $\tau_0 = a_0/\alpha c$ the atomic unit of
time, $\mu = m_p/2m_e = 918.08$, $\omega_\ell = 2\ell_>-1$ the weight of
the lower sublevel, and $D_{ji}$ the dimensionless dipole coupling. The
prefactor $a_0^{3}/\tau_0$ carries the units and everything else is
dimensionless, so $q_{ji}$ is a rate coefficient. Note $D_{ji} \sim 6n^{4}$ for
$\ell_> \ll n$, recovering the $n^{4}$ of Eq.~\eqref{eq:lmix_scaling}, and that
temperature enters as $\Te^{-1/2}$, recovering the $T_i^{-1/2}$ under
$T_i = \Te$. The upward rate follows by reciprocity,
$q_{\mathrm{up}}/q_{\mathrm{down}} =
(2\ell_{\mathrm{hi}}+1)/(2\ell_{\mathrm{lo}}+1)$, which is
Eq.~\eqref{eq:detailed_balance} for degenerate states, where the Boltzmann factor
is unity and only the weights survive.

$F(U_m)$ is the Coulomb logarithm of the physical picture, in the form Badnell's
Eq.~9 gives it:
\begin{equation}
  F(U_m) = \frac{\sqrt{\pi}}{2}U_m^{-3/2}\mathrm{erf}\!\left(\sqrt{U_m}\right)
           - \frac{\mathrm{e}^{-U_m}}{U_m} + E_1(U_m),
  \qquad
  U_m = \frac{E_{\min}}{k_B\Te},
  \label{eq:F_Um}
\end{equation}

\noindent with $E_{\min}$ the minimum energy transfer set by the cut-off radius
$R_c = \lambda_D$. $F$ is dimensionless and \emph{diverges} logarithmically as
$U_m \to 0$: a smaller minimum energy transfer means a larger cut-off radius
means more decades of impact parameter contributing. That behaviour is the
subject of Section~\ref{sec:atoms_corrections}, because an earlier version of
the code had the limit backwards.

Three approximations enter and are carried through the thesis: $T_i = \Te$,
$\nion = \nel$ by quasi-neutrality, and $\Delta\ell = \pm1$ only. A fourth is
easy to miss on reading the code. The Debye length inside $U_m$ is evaluated at
a fixed reference density \SI{e14}{\per\cubic\centi\metre} rather than at each
grid density, on the stated grounds that $F$ depends on density only
logarithmically. That stated ground is wrong, and quoting one number per shell
understates by how much. Evaluating the module's own expression across the eight
density columns at \SI{1}{\electronvolt}, $F$ varies by a factor $3.7$ for
$n=2$, which has only the $\ell=1$ channel, and by between $50$ and $687$ across
the seven $\ell$ channels of $n=8$, the largest belonging to $\ell=1$ and the
smallest to $\ell=7$. An earlier draft quoted $50$ for $n=8$ beside $3.7$ for
$n=2$, which are correct measurements of different $\ell$ channels and should
not have been set side by side. Over the whole $(\Te,\nel)$ grid the $n=8$,
$\ell=1$ spread reaches $1564$
(\texttt{compute\_lmix.py}, functions \texttt{\_psm20\_Um} and
\texttt{\_psm20\_F}). The practice survives for a different reason, measured
rather than assumed. Rebuilding $\Lmat$ with a density-consistent cutoff at every
grid point moves $\ffrac{3}-\ffrac{4}$ by at most \SI{0.42}{\percent}, and
scaling every $\ell$-mixing rate over a range of $\times0.1$ to $\times10$ moves
it by less than \SI{0.3}{\percent} at the benchmark and the density maximum.
$\ell$-mixing is saturated, so the frozen table costs nothing that can be
measured. The
stored table has no density axis at all, so the linear $\propto\nel$ scaling
applied at assembly is exact while the Coulomb-logarithm factor inside it is
frozen at mid-grid. One item is left open rather than settled: Badnell's
formulation offers competing cut-offs to the Debye radius, one set by the level
lifetime and one by any residual splitting of the $\ell$ degeneracy, and the
correct $R_c$ is the smallest of the three. Only the Debye cut-off is
implemented. At these densities it is expected to be the binding one, but that
has not been tested across the grid.

\subsection{Without it, \texorpdfstring{$n=2$}{n=2} is not a physical manifold}

We can now close the loop that Sections~\ref{sec:state_space},
\ref{sec:excitation} and~\ref{sec:radiative} each opened, and the result is the
strongest single statement in this chapter.

Ask what an atom in $2s$ can do. Its only radiative exit is two-photon emission
at \SI{8.229}{\per\second}: $2s \to 1s$ is E1-forbidden, and $\gamma_{2s} = 0$
in this dataset because the two-photon channel is not carried in it. That rate
is ten orders of magnitude below the collisional rates that follow, which is why
its omission is harmless here and why it nonetheless matters: it is the one
depopulation route that does not vanish as $\nel \to 0$. It cannot be moved to $2p$ by an electron: $\Delta n = 0$
transitions were excluded from the CCC dataset on the provider's instruction.
What is left is electron-impact excitation to $n \ge 3$, electron-impact
ionisation, and proton $\ell$-mixing. Decomposing the total loss from $2s$ into
those channels, from the stored coefficient arrays at the benchmark point:
proton $\ell$-mixing $\SI{1.4745e11}{\per\second}$; electron-impact collisions
to $n \ge 3$ $\SI{3.4030e7}{\per\second}$; electron-impact ionisation
$\SI{6.9216e6}{\per\second}$; radiative decay zero. $\ell$-mixing carries
\SI{99.972}{\percent} of everything that happens to a $2s$ atom.

Because $\gamma_{2s}$ is exactly zero, every remaining channel is collisional
and scales with $\nel$, so that share is independent of density: it depends only
on temperature, running from \SI{99.994}{\percent} at
$\Te = \SI{1}{\electronvolt}$ to \SI{99.921}{\percent} at
\SI{10}{\electronvolt}. \textbf{Over the whole 400-point grid, proton
$\ell$-mixing carries at least \SI{99.92}{\percent} of the $2s$ loss rate}, the
minimum falling at the hot edge. There the contrast is starkest: without
$\ell$-mixing $2s$ would empty at \SI{6.17e5}{\per\second}; with it, at
\SI{7.79e8}{\per\second}, a factor of 1262.

The conclusion is not that omitting $\ell$-mixing perturbs the $n=2$ manifold.
It is that omitting it changes which process controls the manifold, by three
orders of magnitude, and leaves $2s$ an artificial metastable trap. Any model of
hydrogen that resolves $\ell$ and excludes $\Delta n = 0$ electron impact
\emph{must} include proton $\ell$-mixing or it is not describing the $n=2$ shell
at all.

\paragraph{Worked numbers.} At the benchmark point
$q_{2s\to2p} = \SI{1.0612e-3}{\cubic\centi\metre\per\second}$
(\texttt{K\_lmix.npy}, element $[2,1,23]$), giving
\SI{1.4745e11}{\per\second} and a $2s$ lifetime of \SI{6.78}{\pico\second}.
Within $n=5$ the $\ell$-mixing loss rates run from \SI{3.25e12}{\per\second} for
$5s$ (\SI{0.31}{\pico\second}) to \SI{3.47e11}{\per\second} for $5g$
(\SI{2.88}{\pico\second}). An independent hand estimate from the
impact-parameter picture, with the Debye length as cut-off and
$b_{\min} = n^{2}a_0$, gives \SI{0.16}{\pico\second} for $n=5$: high by roughly
six, in the expected direction, because the crude estimate uses the full orbital
size for the dipole matrix element and drops the angular factors that $D_{ji}$
supplies. An order-of-magnitude derivation agreeing to within a factor of six
for an understood reason is doing its job: it confirms that the large number
in the table is physics and not a coding error. Finally, the largest single
entry anywhere in the assembled rate matrix at the benchmark point is the $8s$
diagonal, $-\SI{9.5235e12}{\per\second}$ (\texttt{L\_grid.npy}, element
$[23,5,28,28]$): again $\ell$-mixing, at the highest resolved shell, exactly
where the $n^{4}$ of Eq.~\eqref{eq:lmix_scaling} puts it. The stiffest direction
of the whole problem is a process most treatments omit.

%-----------------------------------------------------------------------------
\section{How well do we know these rates?}
\label{sec:atomic_uncertainty}
% QUESTION: where would the atomic data let this model down, and how would we
% know?

Internal consistency checks cannot establish that a rate is the right size.
Detailed balance holds to $2.5\times10^{-9}$ because it was imposed;
column-sum conservation, which Chapter~\ref{ch:theory} uses, holds to
$3.6\times10^{-11}$ because every transition entered as a loss also entered as a
gain. Both would hold for a dataset uniformly wrong by a factor of three. The
only test with real severity is comparison against an independent calculation.

That comparison is against Anderson \emph{et al.}, who computed effective
collision strengths $\Upsilon$ for all $n \le 5$, $\Delta n \ne 0$ transitions in
hydrogen using the R-matrix with pseudostates (RMPS) method, a different
formulation, basis and code from CCC. Two papers carry that calculation and they
are not interchangeable. The original tabulation~\cite{Anderson2000} was
superseded by a corrigendum~\cite{Anderson2002} after Dere and Mason pointed out
that its dipole effective collision strengths did not approach the correct
high-energy asymptotic form. The underlying collision strengths were sound; the
fault was in the Maxwell averaging. The high-energy tail had been fitted by least
squares using the infinite-energy limit point, but RMPS collision strengths
oscillate shallowly at all energies because of high-lying pseudo-thresholds, so
the fit degraded progressively beyond the last computed energy. The corrigendum
recomputes $\Upsilon$ from the same collision strengths using linear
interpolation of the reduced collision strength between the highest finite energy
and the infinite-energy limit. \textbf{This work uses the corrigendum values
throughout.}

How much that choice matters depends entirely on temperature. Comparing the two
tabulations transition by transition, the mean change over $\Te = 1$ to
\SI{10}{\electronvolt} is \SI{1.4}{\percent} and no transition moves by more
than \SI{13.1}{\percent}; at \SI{15}{\electronvolt} the largest change is
\SI{33}{\percent} and at \SI{25}{\electronvolt} it reaches \SI{86}{\percent},
on $1s \to 3p$ in both cases. The revision is therefore almost irrelevant to
this thesis, which never leaves $\Te \le \SI{10}{\electronvolt}$, and decisive
for any extension of this model to attached-divertor or scrape-off-layer
temperatures. The distinction is recorded here because the superseded table
remains in circulation, and because an earlier version of this benchmark used
it.

The comparison itself is done in
\texttt{src/validation/anderson2002\_benchmark.py}, which reads $\Upsilon$
directly from the published corrigendum table rather than from a transcription,
and converts it to a rate coefficient through the authors' Eq.~(3). That
conversion has one trap worth naming: the statistical weight in the denominator
is that of the \emph{lower} state and is $(2S+1)(2L+1) = 2(2\ell+1)$ for
hydrogen doublets, so using $2\ell+1$ introduces a systematic factor of two. The
table's row convention, upper state first, is verified at run time rather than
assumed, by checking that the radiative rate listed against the $2p$--$1s$ row is
$\SI{6.27e8}{\per\second}$. Comparison is made at all 50 temperatures of the
working grid, with $\Upsilon$ interpolated logarithmically in $\Te$; since the
grid spans \SIrange{1}{10}{\electronvolt} and the RMPS table spans
\SIrange{0.5}{25}{\electronvolt}, this interpolates and never extrapolates.
Table~\ref{tab:anderson} gives the outcome over the resulting 4250 comparisons.

\begin{table}[tbp]
  \centering
  \caption[CCC benchmarked against Anderson RMPS]{The CCC excitation rate
  coefficients used in this model, compared against the R-matrix with
  pseudostates calculation of Anderson \emph{et al.}~\cite{Anderson2000}
  as revised by the corrigendum~\cite{Anderson2002}, over 85 matched transitions
  at all 50 temperatures of the working grid (4250 comparisons). The error is
  $(K_{\mathrm{CCC}}/K_{\mathrm{RMPS}} - 1)$; a negative median means CCC is the
  smaller of the two. Source:
  \texttt{data/processed/collisions/ccc\_vs\_anderson2002\_thesis\_Te\_grid.csv}
  and the report \texttt{ccc\_anderson.md}. The agreement degrades smoothly with
  the upper shell rather than failing abruptly at $n=5$.}
  \label{tab:anderson}
  \begin{tabular}{lrrrr}
    \toprule
    Subset & comparisons & within \SI{20}{\percent} & mean $|$error$|$ & median error \\
    \midrule
    $n_{\text{upper}} = 2$     &  100 & \SI{100.0}{\percent} & \SI{10.9}{\percent} & \SI{-10.5}{\percent} \\
    $n_{\text{upper}} = 3$     &  450 & \SI{96.9}{\percent}  & \SI{8.5}{\percent}  & \SI{-5.6}{\percent} \\
    $n_{\text{upper}} = 4$     & 1200 & \SI{74.0}{\percent}  & \SI{14.9}{\percent} & \SI{-9.9}{\percent} \\
    $n_{\text{upper}} = 5$     & 2500 & \SI{12.3}{\percent}  & \SI{41.1}{\percent} & \SI{-34.2}{\percent} \\
    \midrule
    $n_{\text{upper}} \le 4$   & 1750 & \SI{81.4}{\percent}  & \SI{13.0}{\percent} & \SI{-8.8}{\percent} \\
    all transitions            & 4250 & \SI{40.7}{\percent}  & \SI{29.5}{\percent} & \SI{-21.5}{\percent} \\
    \bottomrule
  \end{tabular}
\end{table}

\begin{figure}[tbp]
  \centering
  \includegraphics[width=\textwidth]{fig2_2_ccc_rmps_benchmark}
  \caption[CCC benchmarked against RMPS]{\CapFigCCCRMPS}
  \label{fig:ccc_rmps}
\end{figure}

Taken at face value the last row is a failure: fewer than half the comparisons
agree within \SI{20}{\percent}. Most of that sits in transitions reaching $n=5$,
which supply 2500 of the 4250 comparisons. But the degradation with upper shell
is graded rather than abrupt, and saying that the disagreement lies
\emph{entirely} at $n=5$ would overstate the case: $n_{\text{upper}} = 4$ already
falls to \SI{74.0}{\percent} within \SI{20}{\percent} against
\SI{96.9}{\percent} for $n_{\text{upper}} = 3$. Within the $n_{\text{upper}}=4$
block the spread is carried by $\Delta n = 1$ and $\Delta n = 3$; the $\Delta n =
2$ transitions agree within \SI{20}{\percent} at every point tested. Nor is the
agreement uniform in temperature. For $n_{\text{upper}} \le 4$ the mean absolute
deviation runs from \SI{7.2}{\percent} at \SI{1}{\electronvolt} to
\SI{17.9}{\percent} at \SI{10}{\electronvolt}, so the figure of
\SI{13}{\percent} in Table~\ref{tab:anderson} is a grid average and not a
bound. Figure~\ref{fig:ccc_rmps} shows both features: the $n'=5$ points
scatter over more than a factor of ten while the $n' \le 3$ points stay inside
the \SI{20}{\percent} band, and the size of the scatter is independent of the
size of the rate coefficient (rank correlation of $|$error$|$ with $\log K$ of
$0.02$ over the $2500$ comparisons at $n'=5$), which rules out a near-threshold
or interpolation artefact in the Maxwell averaging as the cause. Its sign is
not: CCC reads low at the smallest rates and high at the largest, which are
the two modes described next.

At $n=5$ the disagreement has two modes with opposite signs. Where the
transition is dipole-allowed between the nearly degenerate $n=4$ and $n=5$
shells, as for $4s \to 5p$ with its gap of \SI{0.31}{\electronvolt}, CCC gives
the larger rate, by up to a factor of $2.97$; every one of the $105$
comparisons in which CCC exceeds RMPS by more than a factor of two belongs to
the four transitions $4s \to 5p$, $4p \to 5s$, $4p \to 5d$ and $4d \to 5f$,
and the excess grows with temperature, from $1.4$ at \SI{1}{\electronvolt} to
$3.0$ at \SI{5}{\electronvolt} for $4s \to 5p$. Where the transition is a
high-multipole one with no dipole matrix element, as for $1s \to 5g$ in which
$\ell$ changes by four, CCC gives rates several times smaller: a factor of $5$
for $1s \to 5g$ at every temperature on the grid, and $275$ comparisons across
eleven transitions fall below half the RMPS value. This second mode dominates
the aggregate, which is why the median error at $n_{\text{upper}}=5$ is
\SI{-34.2}{\percent}. The dipole transitions out of the ground state sit
between the two, with CCC low by a factor that shrinks with temperature,
$0.54$ at \SI{1}{\electronvolt} to $0.75$ at \SI{10}{\electronvolt} for
$1s \to 5p$. An earlier draft counted ten comparisons beyond a factor of two,
all with CCC the larger; that was a count over four temperatures of one
direction only, and the figure plots all fifty.

\subsection{Which calculation is anomalous at \texorpdfstring{$n=5$}{n=5}}
\label{sec:n5_anomaly}

Both calculations are approximate and neither is a measurement, so the
disagreement cannot be settled by preferring one author. It can, however, be
\emph{localised}, using an asymmetry between the two datasets: the RMPS
tabulation stops at $n=5$, while CCC continues to $n=10$. Only one of them can
be asked what happens next.

The test uses a property of hydrogen that holds regardless of which code
computes it. Along a Rydberg series of dipole transitions $n_{\mathrm{lo}} \to
n_{\mathrm{up}}$ the absorption oscillator strength falls as
$f \sim C/n_{\mathrm{up}}^{3}$ as $n_{\mathrm{up}} \to \infty$, and at a
temperature well above threshold the rate coefficient is dominated by the Bethe
term --- the leading high-energy contribution to a dipole-allowed cross section,
which is proportional to $f$ --- and inherits that scaling. The rate
coefficient also carries the threshold factor $\mathrm{e}^{-\Delta E/\Te}$,
with $\Delta E$ the excitation energy of the transition, which both
calculations share and which at \SI{1}{\electronvolt} dominates everything
else: along $1s \to np$ it alone falls by $26$, $15$ and \SI{10}{\percent} at
the steps into $n = 5$, $6$ and $7$. Dividing it out leaves the reduced
quantity
\begin{equation}
  R(n_{\mathrm{up}}) = K(n_{\mathrm{lo}} \to n_{\mathrm{up}})\,
  n_{\mathrm{up}}^{3}\, \mathrm{e}^{\Delta E/\Te},
  \label{eq:rydberg_reduced}
\end{equation}
proportional to $\Upsilon\,n_{\mathrm{up}}^{3}$ over the weight of the lower
state, which must flatten as the series is climbed.
Equation~\eqref{eq:rydberg_reduced} does not say what the constant is, and it
is not used to compute anything; its only role is that a dataset tracking
hydrogenic behaviour produces a smooth $R(n_{\mathrm{up}})$, so a dataset that
kinks at one particular $n_{\mathrm{up}}$ is departing from that behaviour
there. An earlier draft tabulated $K n_{\mathrm{up}}^{3}$ with the threshold
factor left in; at \SI{1}{\electronvolt} that quantity is smooth for either
dataset because the Boltzmann factor is, and the $-41$, $-16$ and
\SI{-12}{\percent} it recorded for CCC were the threshold factor rather than
the scaling under test. Table~\ref{tab:rydberg} gives the successive changes
in $R$ for three series at both ends of the temperature range.

\begin{table}[tbp]
  \centering
  \caption[Rydberg-series test of the $n=5$ disagreement]{Successive percentage
  changes in the reduced rate coefficient $R(n_{\mathrm{up}})$ of
  Eq.~\eqref{eq:rydberg_reduced}, threshold factor removed, along three dipole
  series at \SI{1}{\electronvolt} and \SI{10}{\electronvolt}. At
  \SI{1}{\electronvolt} the RMPS step into $n=5$, the top shell of their
  target basis, is a \emph{rise} in all three series, the only positive
  entries in the table; at \SI{10}{\electronvolt} neither calculation kinks
  and the RMPS steps are uniformly the flatter. Dashes mark shells the RMPS
  tabulation does not contain. Source:
  \texttt{data/processed/collisions/anderson\_validity\_dipole\_series.csv},
  columns \texttt{R\_CCC\_bz} and \texttt{R\_And\_bz}.}
  \label{tab:rydberg}
  \begin{tabular}{lllrrrrr}
    \toprule
    Series & $\Te$ & Source & $2\!\to\!3$ & $3\!\to\!4$ & $4\!\to\!5$ & $5\!\to\!6$ & $6\!\to\!7$ \\
    \midrule
    $1s \to np$ & \SI{1}{\electronvolt}  & CCC  & $-19.5$ & $-9.2$  & $-20.6$          & $-1.2$  & $-2.5$ \\
                &                        & RMPS & $-19.6$ & $-9.9$  & $\mathbf{+28.4}$ & ---     & --- \\
                & \SI{10}{\electronvolt} & CCC  & $-41.4$ & $-17.9$ & $-9.9$           & $-5.2$  & $-3.8$ \\
                &                        & RMPS & $-34.0$ & $-10.9$ & $-6.9$           & ---     & --- \\
    \midrule
    $2p \to nd$ & \SI{1}{\electronvolt}  & CCC  & ---     & $-42.6$ & $-27.3$          & $-14.9$ & $-8.3$ \\
                &                        & RMPS & ---     & $-42.6$ & $\mathbf{+12.4}$ & ---     & --- \\
                & \SI{10}{\electronvolt} & CCC  & ---     & $-64.6$ & $-33.2$          & $-18.5$ & $-11.8$ \\
                &                        & RMPS & ---     & $-59.1$ & $-29.3$          & ---     & --- \\
    \midrule
    $2s \to np$ & \SI{1}{\electronvolt}  & CCC  & ---     & $-28.5$ & $-24.2$          & $-10.2$ & $-10.1$ \\
                &                        & RMPS & ---     & $-29.8$ & $\mathbf{+12.0}$ & ---     & --- \\
                & \SI{10}{\electronvolt} & CCC  & ---     & $-54.4$ & $-26.6$          & $-14.4$ & $-9.6$ \\
                &                        & RMPS & ---     & $-47.1$ & $-21.5$          & ---     & --- \\
    \bottomrule
  \end{tabular}
\end{table}

\begin{figure}[tbp]
  \centering
  \includegraphics[width=\textwidth]{fig2_3_rydberg_series}
  \caption[Which calculation is anomalous at $n=5$]{\CapFigRydberg}
  \label{fig:rydberg}
\end{figure}

Figure~\ref{fig:rydberg} plots the same quantity against shell at both ends
of the temperature range.

The two datasets agree on the shape of each series until the last step the
RMPS calculation can take, and what happens at that step depends on the
temperature. At \SI{1}{\electronvolt} the RMPS reduced rate \emph{rises} into
$n=5$, by $28$, $12$ and \SI{12}{\percent} in the three series, while every
other step of either calculation in Table~\ref{tab:rydberg} is a fall; the
ratio of the RMPS rate to the CCC rate along $1s \to np$ is $1.15$, $1.15$,
$1.14$ and then $1.85$. CCC is not smooth there either, dropping by
\SI{20.6}{\percent} into $n=5$ on $1s \to np$ between neighbouring steps of
$-9$ and \SI{-1}{\percent}, so both calculations carry structure at that shell
and the localisation rests on the sign: a rise in
$\Upsilon\,n_{\mathrm{up}}^{3}$ along a Rydberg series is what a truncated
basis produces, and a smooth continuation is not. At \SI{10}{\electronvolt}
there is no kink. The RMPS steps are three to seven points flatter than the
CCC ones at $3 \to 4$ and $4 \to 5$ alike, the step into $n=5$ is where the
two calculations differ \emph{least}, and the same ratio climbs steadily,
$1.06$, $1.19$, $1.29$, $1.34$. An earlier draft placed the kink at both
temperatures; it is a \SI{1}{\electronvolt} feature, and at
\SI{10}{\electronvolt} the disagreement is a trend with $n$ rather than a
step at one shell.

The sign at \SI{1}{\electronvolt} matches the bulk sign of the disagreement
in Table~\ref{tab:anderson}: the RMPS $n=5$ rates sit high, so CCC reads low
against them. The temperature dependence constrains the mechanism. A top
shell that absorbed flux belonging to $n \ge 6$ and the continuum would do so
at every collision energy, and the excess would appear at every temperature;
it appears only where the Maxwell average samples the cross section within
about an electron-volt of the $n=5$ threshold, which is where the $n \ge 6$
thresholds, the ionisation limit and the pseudo-state thresholds of a
truncated expansion crowd together. The evidence points to a near-threshold
artefact of the basis rather than a high-energy one.

This localises the anomaly; it does not prove CCC correct. The scaling
argument is asymptotic, and \SI{1}{\electronvolt} is not a temperature at
which the Bethe term cleanly dominates, so the test is suggestive rather than
decisive. What it establishes is narrower and still useful: the departure
from hydrogenic behaviour occurs at one specific shell of one specific
calculation, at the temperatures where the two disagree most, rather than
being distributed through the CCC dataset. \textbf{This work therefore treats
the $n=5$ disagreement as a limitation of the benchmark rather than as a
measured error in the data being benchmarked}, and claims no more than that.

\subsection{Does any of it reach the result?}
\label{sec:atomic_propagation}

The preceding argument is about rate coefficients. A discrepancy in a rate
coefficient matters only insofar as it moves something the thesis reports, and
the way to find out is to make the substitution and look. Because the CR matrix
is linear in electron density --- verified, not assumed, to a relative
$3\times10^{-16}$ --- it separates exactly into a radiative part and a
collisional part,
\begin{equation}
  \Lmat(\Te, \nel) = \Lmat_{\mathrm{rad}}(\Te) + \nel\,\Lmat_{\mathrm{coll}}(\Te),
  \label{eq:L_split}
\end{equation}
so the collisional rate coefficients can be replaced without disturbing
anything else. For all 85 transitions the RMPS calculation covers, the CCC rate
coefficient in $\Lmat_{\mathrm{coll}}$ was replaced by the RMPS one and the eigenvalue problem
of Chapter~\ref{ch:theory} re-solved at every one of the 400 grid points.
Excitation and de-excitation were scaled by the same factor, so detailed balance
survives the substitution, and the diagonal was compensated so that column sums,
and with them the loss to the continuum, are unchanged. Before any substitution
the procedure reproduces the unperturbed spectrum, including all three ranges of
Eq.~\eqref{eq:M_range}, which is the check that makes the perturbed numbers
worth reading.

Replacing every rate coefficient the benchmark covers changes $\tauslow$ by
\SI{11.4}{\percent} on average across the grid and by at most
\SI{18.7}{\percent}; $\taurel$ moves by \SI{2.4}{\percent}; and $\Msep$, which
inherits both, moves by \SI{11.1}{\percent} on average. At the benchmark point
$\Msep$ falls from 9982 to 8855. This is the test that
Section~\ref{sec:n5_anomaly} could not perform and that an earlier draft of this
chapter identified as missing.

\begin{figure}[tbp]
  \centering
  \includegraphics[width=\textwidth]{fig2_4_atomic_propagation}
  \caption[What the atomic-data disagreement costs]{\CapFigAtomicPropagation}
  \label{fig:atomic_propagation}
\end{figure}

Figure~\ref{fig:atomic_propagation}(a) maps that shift over the grid. It grows
with density and saturates above $\nel \approx
\SI{e13}{\per\cubic\centi\metre}$, and it is largest at the cold end --- the
opposite of where the rate coefficients are best known, because at
\SI{1}{\electronvolt} the excitation rates sit deep in the Boltzmann
exponential and a given fractional change in a rate produces a larger one in
the eigenvalue it feeds.

Its outcome is not the expected one, and the useful part is the disagreement
with what was expected. Decomposing the substitution by which transitions are
replaced:

\begin{table}[tbp]
  \centering
  \caption[Which transitions carry the atomic-data sensitivity]{Change in the
  timescale separation $\Msep$ when the CCC rate coefficients are replaced by
  the RMPS values, for subsets of the 85 benchmarked transitions. The
  50~transitions reaching $n=5$, which carry almost all of the benchmark
  disagreement, move $\Msep$ by under \SI{1}{\percent}; the 14~ground-state
  transitions, which carry much less of it, move $\Msep$ by more than the full
  substitution does. Source:
  \texttt{validation/anderson\_subset\_impact/}, which repeats the
  substitution of \texttt{ccc\_anderson\_grid\_impact.py} for each subset.}
  \label{tab:atomic_sensitivity}
  \begin{tabular}{lrrr}
    \toprule
    Transitions replaced & count & mean $|\Delta\tauslow|$ & $\Delta\Msep$ at benchmark \\
    \midrule
    all benchmarked ($n_{\text{upper}} \le 5$) & 85 & \SI{11.4}{\percent} & \SI{-11.3}{\percent} \\
    $n_{\text{upper}} = 5$ only                & 50 & \SI{4.7}{\percent}  & \SI{-0.9}{\percent} \\
    $n_{\text{upper}} \le 4$ only              & 35 & \SI{7.2}{\percent}  & \SI{-10.8}{\percent} \\
    ground-state $1s \to n\ell$ only           & 14 & \SI{10.4}{\percent} & \SI{-12.4}{\percent} \\
    \bottomrule
  \end{tabular}
\end{table}

Fifty transitions reaching $n=5$, disagreeing by factors of two to five, move
$\Msep$ by \SI{0.9}{\percent}. Fourteen ground-state transitions, agreeing far
better, move it by \SI{12.4}{\percent}. The reason is structural rather than
numerical: $\tauslow$ is the timescale of the slowest mode, and
Section~\ref{sec:two_clocks} identifies that mode as the depletion of the ground
state, so the rates out of $1s$ set it almost by themselves. Rates into $n=5$
enter only through cascades that have already been diluted by everything they
pass through. The \SI{0.9}{\percent} is also a net: the five $1s \to 5\ell$
transitions alone move $\Msep$ by \SI{-3.3}{\percent} and the forty-five
excited-fed transitions into $n=5$ by \SI{+2.4}{\percent}, so part of the
reason $n=5$ does not reach the result is that its two halves pull in opposite
directions (\texttt{validation/anderson\_subset\_impact/}).

The conclusion of the earlier draft survives and its reasoning does not. It is
true that the $n=5$ disagreement does not affect this thesis, but not because
$n=5$ channels are weak in some general sense --- it is because they are weak
\emph{in the specific eigenvalue the thesis reports}, and that had to be
computed rather than asserted. The same calculation shows that the transitions
which do matter are the ones the benchmark constrains least well among
$n_{\text{upper}} \le 4$: ground-state excitation agrees within
\SI{20}{\percent} in only \SI{61.8}{\percent} of comparisons at
$n_{\text{upper}} \le 4$, against \SI{81.4}{\percent} for that block as a whole.
Chapter~\ref{ch:theory} carries the consequence, which is that $\Msep$ is
reported with an atomic-data systematic of \SI{11}{\percent}.

Three further soft spots are worth naming rather than leaving to be discovered.
\emph{Detailed balance in the raw cross sections}: separately from the enforced
consistency above, the CCC cross sections were checked for microscopic
reversibility as supplied (\texttt{src/parsers/qc\_ccc.py}) and satisfy it at
the \SI{0.05}{\percent} level. That is a genuine physics check on the provider's
calculation, and it passes comfortably. \emph{Lotz ionisation for the top six
shells, and its echo in recombination}: the metadata records that Lotz
``overestimates CCC by factor \textasciitilde4--8'' for $n = 10$--$15$, and
because Eq.~\eqref{eq:alpha_3BR} builds $\alpha_{\mathrm{3BR}}$ \emph{from} $Q$,
that overestimate propagates straight into the three-body coefficients of the
same shells, the shells whose feed is essentially all three-body.
Chapter~\ref{ch:limits} reports, from an independent comparison against
published population coefficients, that the terminal shell is over-populated by
4.9--6.2 times while $n=10$ agrees to \SIrange{4}{20}{\percent}. The two
observations agree in size and location, but the test that would establish the
link has not been done, so the connection is a hypothesis, not a finding.
\emph{Vriens--Smeets for the bundled block}: 201 of the 819 excitation pairs
come from a semi-empirical fit of quoted accuracy $\sim$\SI{20}{\percent} rather
than a close-coupling calculation, which is the price of the ladder reaching
$n=15$. And one gap that the benchmark simply cannot close: the RMPS tabulation
stops at $n=5$, so every CCC transition reaching $n \ge 6$ --- the whole upper
half of the ladder --- has no independent check in this thesis at all.

%-----------------------------------------------------------------------------
\section{Two errors found by audit}
\label{sec:atoms_corrections}
% QUESTION: what did we get wrong while building this, how was it found, and
% how large was it?

Two implementation errors were found in this pipeline by auditing it against its
own published sources. Neither was found by a test failing; both were found by
reading the code against the paper it claimed to implement, and by recomputing a
quantity unconditionally and comparing with the stored value. They belong in the
body of the thesis for two reasons: a reader assessing the numbers is entitled
to know what the audit turned up, and the \emph{size} of each error, once
bounded, is a physical result in its own right.

\subsection{The \texorpdfstring{$\ell$}{l}-mixing Coulomb logarithm}

An earlier version of \texttt{src/rates/compute\_lmix.py} set the factor $F(U_m)$
of Eq.~\eqref{eq:F_Um} to $1$, justified in a comment as ``the low-density limit
$U_m \to 0$, $F \to 1$''.

The justification inverts the limit. Reading Badnell's Eq.~9~\cite{Badnell2021}
directly from the source, $F(U_m)$ \emph{is} the Coulomb logarithm of the
impact-parameter picture, and $F(U_m \to 0)$ diverges logarithmically rather than
tending to one, as it must, since a vanishing minimum energy transfer means an
unbounded cut-off radius. The values of $U_m$ at divertor conditions are indeed
small, $10^{-3}$ to $10^{-1}$, which is precisely where $F$ is \emph{large}. At
the benchmark temperature the correct factor for $2p \to 2s$ is $F = 6.95$ at
$U_m = 1.05\times10^{-3}$; across transitions it runs from about 7 at low $n$ to
about 3 at $n=8$. The old code under-counted every $\ell$-mixing rate by 3 to 7.

That the correction is present in the data this thesis uses can be checked
rather than taken on trust. The canonical matrix is
\texttt{cr\_matrix/L\_grid.npy}, written 21 July 2026, SHA-256 beginning
\texttt{2d92b58e1693}; the $\ell$-mixing coefficient it was built from is
$q_{2p\to2s} = \SI{3.53722e-4}{\cubic\centi\metre\per\second}$ at temperature
index 23, while evaluating Eq.~\eqref{eq:psm20} independently with $F$ set to 1
gives \SI{5.08977e-5}{\cubic\centi\metre\per\second}. The ratio is $6.9497$ and
the independently computed $F(U_m)$ at the same point is $6.94961$, so the
stored coefficient equals $q\,F$ to 8 parts in $10^{6}$.

What the correction changed is the interesting part. Scaling every intra-shell
off-diagonal by the correct $F$, restoring the exact column sums and recomputing
the spectrum moves the excited-state relaxation time by \SI{-0.12}{\percent} and
the ionisation-balance time by \SI{+0.07}{\percent} at the benchmark point;
across the grid the effect on the relaxation time is bounded by
\SI{0.85}{\percent} (register entry A12). A factor of 3 to 7 in the fastest
rates in the matrix moves the slow eigenvalues by less than one percent.

That insensitivity is not luck and it is not an excuse. It is the timescale
separation of Chapter~\ref{ch:theory} doing what timescale separation does.
$\ell$-mixing runs at picoseconds; the relaxation being measured runs at
nanoseconds, and on that timescale the $\ell$ sublevels are fully equilibrated
whether the mixing is three times faster or seven, so the slow eigenvalue cannot
see the difference. Finding a factor-of-five error in an input and then proving
the headline moves by less than one percent is a stronger statement about the
model than never having looked. It licenses nothing elsewhere: absolute
$\ell$-resolved populations and any $\ell$-resolved line ratio depend on the
correction directly, which is why it was fixed rather than merely bounded.

\subsection{The eigenvalue magnitude filter}

Two analysis scripts extracted timescales by sorting the eigenvalues of the rate
matrix and then discarding every eigenvalue of magnitude below
\SI{1}{\per\second}, with the literal \texttt{eigs = eigs[eigs < -1.0]} and a
comment calling the discarded values ``near-zero numerical noise''.

The slowest genuine eigenvalue of this system is not noise. It is the
ionisation-balance rate, and in the cold, thin corner of the grid its magnitude
falls below \SI{1}{\per\second}: the corresponding timescale reaches tens of
seconds. Wherever that happened the filter deleted the physical slow mode and
promoted the whole eigenvalue ladder by one rung, so the reported
ionisation-balance time was really the relaxation time of the level above it.

19 of the 400 grid points were affected, all at $\Te \le
\SI{1.389}{\electronvolt}$ and $\nel \in
[10^{12}, 1.93\times10^{13}]\,\si{\per\cubic\centi\metre}$. At every one the
corrupted relaxation time equalled the corrupted ionisation-balance time of the
unfiltered calculation exactly, the signature of a one-rung shift, and how the
mechanism was confirmed rather than guessed. The reported range of the
timescale-separation ratio changed from $1.34$--$1.01\times10^{8}$ to
$86.77$--$1.72928\times10^{9}$. The benchmark point was bit-identical before and
after, and so were the step-error maps at all 400 points: the corruption was
real, confined, and did not touch the central result. Every check the filtered
code ran continued to pass, so it was found not by a check failing but by
recomputing the spectrum unconditionally, storing it separately, and comparing
bit-for-bit. The filtered outputs were kept rather than deleted, under names
ending \texttt{\_FILTERED\_20260721}, so the comparison remains reproducible.

Honesty requires the last statement to be precise rather than tidy. The filter
was first removed from the analysis scripts that write the timescale grids, which
select on \texttt{eigs < 0.0}. The literal \texttt{eigs[eigs < -1.0]} survived
for some weeks afterwards in two other files, \texttt{src/rates/solve\_cr.py}
line 269 and \texttt{src/rates/check\_mz.py} line 10, neither of which produces
any number quoted in this thesis: the first is a superseded solver and the second
a diagnostic print. Both were corrected on 16 September 2026, once it had been
established that nothing imported either, and the literal no longer appears
anywhere in \texttt{src/}. The gap between the two repairs is recorded rather
than smoothed over, because during it a reader who searched the repository for
the literal would have found it and been entitled to ask which numbers passed
through it.

\subsection{What both errors have in common}

Both were invisible to every consistency test the pipeline ran. The
$\ell$-mixing error preserved column sums exactly, because $\ell$-mixing is a
conservative redistribution and its magnitude cancels in the sum. The eigenvalue
filter produced eigenvalues that were all negative, all finite and correctly
ordered. A structurally consistent, physically wrong result is the
characteristic failure mode of this kind of model, and it is why
Chapter~\ref{ch:validation} separates checks a wrong result would also have
passed from checks it would not. The structural response is the one described in
Section~\ref{sec:state_space}: grids, state ordering and matrices are loaded
from the pipeline's own files through a single module that raises on any
inconsistency, so a stale array cannot be silently combined with a fresh one.

%-----------------------------------------------------------------------------
\section{What this model assumes}
\label{sec:atoms_scope}
% QUESTION: what has been left out of the atomic description, and where would
% each omission bite?

Every item here is a deliberate choice, gathered so that they can be found in
one place and quantified where the thesis is able to quantify them.

\begin{description}
  \item[Maxwellian electrons.] Every rate coefficient is an average over a
    Maxwell--Boltzmann distribution, and Eq.~\eqref{eq:detailed_balance}, on
    which every de-excitation and every three-body coefficient rests, holds
    exactly for a Maxwellian and not otherwise. A suprathermal tail, which a
    real transient can produce, would break it.
  \item[$T_i = \Te$ and $\nion = \nel$.] The proton temperature in the
    $\ell$-mixing rate is set to the electron temperature; since
    $\nu_{\ell\text{-mix}} \propto T_i^{-1/2}$, a factor of two error in $T_i$
    is \SI{40}{\percent} in the rate, inside the insensitivity of
    Section~\ref{sec:atoms_corrections}, but stated rather than assumed away.
    Quasi-neutrality gives $\nion = \nel$; no impurities are modelled.
  \item[A fixed reference density in the Debye cut-off.] $U_m$ in
    Eq.~\eqref{eq:F_Um} is evaluated at \SI{e14}{\per\cubic\centi\metre} rather
    than at each grid density; the dependence is logarithmic, and the stated
    variation across the grid is about \SI{30}{\percent} in $F$.
  \item[Optically thin.] All Einstein coefficients enter bare,
    $\Theta_P \equiv 1$. Chapter~\ref{ch:theory} evaluates
    Eq.~\eqref{eq:escape} across the grid and finds the assumption excellent
    over the hot half and failing along the $\Te = \SI{1}{\electronvolt}$ edge;
    Chapter~\ref{ch:limits} quantifies the consequence.
  \item[$n_{\max} = 15$, with $n \ge 9$ bundled.] Seven shells are one state
    each, with no $\ell$-mixing inside the bundled block: they are \emph{assumed}
    in statistical $\ell$-equilibrium rather than driven to it, and that has not
    been tested directly. Chapter~\ref{ch:limits} reports the measured excess in
    the terminal shell.
  \item[Uniform plasma, no transport, no molecules.] A single well-mixed volume
    of atomic hydrogen only, in reactor terms a perfectly stirred tank rather
    than a plug-flow reactor. At the cold end of the grid the gas is largely neutral
    and in reality substantially molecular, and molecular-assisted recombination
    would feed $n=3$ directly. That is a physics limit rather than a data limit,
    and it is why the grid stops at \SI{1}{\electronvolt}.
  \item[An open system.] The ion H$^{+}$ is a fixed external reservoir: there is
    no 44th state, the model does not conserve nuclei, and ionisation is a
    one-way leak out of the tracked manifold. This is why the ionisation-balance
    timescale is boundary-dependent, and why the effective coefficients this
    model produces cannot be compared directly with those tabulated by
    ADAS. No ADAS dataset enters the rate matrix at any point; ADAS appears in
    this thesis only as an external comparison.
\end{description}

%-----------------------------------------------------------------------------
\section{Summary}
\label{sec:atoms_summary}
% QUESTION: what has this chapter established, and what is now possible that
% was not possible before it?

This chapter has built the input side of a collisional--radiative model of
atomic hydrogen: 43 levels, six elementary processes, and a rate coefficient for
every pair of levels at every point of a $50\times8$ grid in temperature and
density. Table~\ref{tab:data_sources} is the provenance record for all of it.

One provenance item on that table was resolved by test rather than by
inspection, and the outcome is recorded here because the suspicion was
reasonable and wrong. The merged excitation arrays appeared to predate a CCC
file they are assembled from, which would have meant
\texttt{assemble\_K\_exc.py} had not been re-run since. Re-running the merge
into a scratch directory and comparing byte for byte returns all five arrays
\emph{identical} by SHA-256, with a maximum absolute difference of zero. The
file whose date raised the question is \texttt{K\_CCC\_metadata.csv}, which
carries no rate data; the rate tables themselves predate the merge. The stored
arrays are current.

\begin{table}[tbp]
  \centering
  \caption[Atomic data sources]{Every atomic dataset entering the model, with
  the file it is stored in and the script that produced it. Paths are relative
  to \texttt{data/processed/}; the arrays are those named in the \texttt{PATHS}
  dictionary of \texttt{src/rates/assemble\_cr\_matrix.py}. All rate-coefficient
  arrays are tabulated on the 50-point temperature grid
  $\Te = 1$--\SI{10}{\electronvolt}. Two hydrogen energies appear in the
  pipeline and they are two \emph{different physical quantities}, not two
  estimates of one; Section~\ref{sec:two_energies} says which is used where.}
  \label{tab:data_sources}
  \small
  \begin{tabular}{p{0.20\textwidth}p{0.25\textwidth}p{0.23\textwidth}p{0.22\textwidth}}
    \toprule
    Quantity & Source & Stored file & Producing script \\
    \midrule
    Excitation and de-excitation, $n \le 10$ &
      CCC (Bray, priv.\ comm.\ 2026); $\Delta n = 0$ excluded &
      \texttt{collisions/K\_exc\_full/K\_exc\_full.npy}, \texttt{K\_deexc\_full.npy} &
      \texttt{compute\_K\_CCC.py}, \texttt{assemble\_K\_exc.py} \\
    Excitation, bundled shells &
      Vriens--Smeets Eqs.~(17)--(19) &
      merged into the same arrays &
      \texttt{compute\_K\_VS.py} \\
    Ionisation, $n \le 9$ &
      CCC total ionisation cross sections &
      \texttt{collisions/tics/K\_ion\_final.npy} &
      \texttt{compute\_K\_TICS.py} \\
    Ionisation, $n = 10$--$15$ &
      Lotz semi-empirical formula &
      same array &
      \texttt{ionization\_rates.py} \\
    Radiative recombination &
      Johnson three-term Gaunt expansion; $\ell$-split statistically &
      \texttt{recombination/alpha\_RR\_resolved.npy}, \texttt{\_bundled.npy} &
      \texttt{recombination\_rates.py} \\
    Three-body recombination &
      detailed balance from ionisation, via Saha &
      \texttt{recombination/alpha\_3BR\_resolved.npy}, \texttt{\_bundled.npy} &
      \texttt{recombination\_rates.py} \\
    Einstein $A$ coefficients &
      Hoang-Binh, ADUU v1.0 &
      \texttt{Radiative/A\_resolved.npy}, \texttt{A\_bund\_res.npy}, \texttt{A\_bund\_bund.npy}, \texttt{gamma\_*.npy} &
      \texttt{radiative\_rates.py} \\
    Proton-impact $\ell$-mixing &
      PSM20 Debye-cutoff formula, Badnell \emph{et al.} &
      \texttt{lmix/K\_lmix.npy} &
      \texttt{compute\_lmix.py} \\
    State ordering and grids &
      the pipeline's own output &
      \texttt{collisions/K\_exc\_full/state\_index.csv}, \texttt{cr\_matrix/Te\_grid\_L.npy}, \texttt{ne\_grid\_L.npy} &
      \texttt{assemble\_cr\_matrix.py}, read via \texttt{cr\_context.py} \\
    \bottomrule
  \end{tabular}
\end{table}

Four things carry forward. The \textbf{state space is not arbitrary}: $\ell$ is
resolved below $n=9$ because $2s$ and $2p$ behave completely differently, and
bundled above it because $\ell$-mixing makes the sublevels of a high shell
indistinguishable, which was checked, not asserted. \textbf{Proton-impact
$\ell$-mixing is structural}: it carries at least \SI{99.92}{\percent} of the
loss from $2s$ everywhere on the grid. \textbf{The data are good where it
matters}: the $n_{\text{upper}} \le 4$ transitions that drive the kinetics agree
with an independent RMPS calculation at a mean absolute error of
\SI{13.0}{\percent}. The agreement degrades smoothly with the upper shell
rather than holding and then failing, and the factor 2--5 disagreement at $n=5$
is localised to the top shell of the RMPS target basis by a test those data
cannot perform on themselves. Substituting the RMPS rates wholesale moves the
timescale separation of Chapter~\ref{ch:theory} by \SI{11}{\percent}, carried
by ground-state excitation and not by $n=5$. And \textbf{two
implementation errors were found by audit and both were bounded}.

What none of this has done is put the atoms in motion. We have $43 \times 43$
rate coefficients and a feed term, and no idea what the system does with them.
Chapter~\ref{ch:theory} assembles them into a single matrix equation, finds that
the network runs on two clocks separated by four orders of magnitude, and asks
the question this thesis was built to answer: whether the fast one is fast
enough for the standard shortcut to be safe.
